\documentclass{article}
\usepackage{amsmath, amsthm, amsfonts, amssymb}
\theoremstyle{plain}
\newtheorem{thm}{Theorem}
\numberwithin{thm}{section}

\theoremstyle{plain}
\newtheorem{prop}{Proposition}
\numberwithin{prop}{section}

\theoremstyle{definition}
\newtheorem{defn}{Definition}
\numberwithin{defn}{section}

\theoremstyle{remark}
\newtheorem*{rem}{Remark}

\newtheorem*{cor}{Corollary}

\numberwithin{equation}{section}

\begin{document}
\title{Sequences and Series}
\author{Amey Joshi}
\maketitle
\section{Sequences of real numbers}\label{s1}
\begin{defn}\label{s1d1}
A sequence is a mapping $\mathbb{N} \mapsto \mathbb{R}$.
\end{defn}

\begin{defn}\label{s1d2}
A sequence $\{x_n\}$ is said to converge to a limit $L$ if for a given 
$\epsilon > 0$, we can find $N \in \mathbb{N}$ such that for all $n > N$, 
$|x_n - L| < \epsilon$.
\end{defn}

\begin{defn}\label{s1d3}
A sequence $\{x_n\}$ is said to be bounded from above (below) if there exists
$M$ ($m$) such that $x_n \le M$ ($x_n \ge m$) for all $n \in N$.
\end{defn}

\begin{defn}\label{s1d4}
A sequence $\{x_n\}$ is said to be bounded if it is bounded from above and below.
\end{defn}

\begin{defn}\label{s1d5}
A sequence $\{x_n\}$ is said to be monotone increasing if $x_n \le x_{n+1}$ for
all $n \in \mathbb{N}$. It is said to be strictly monotone increasing if $x_n <
x_{n+1}$ for all $n \in \mathbb{N}$.
\end{defn}

\begin{defn}\label{s1d6}
A sequence $\{x_n\}$ is said to be monotone decreasing if $x_n \ge x_{n+1}$ for
all $n \in \mathbb{N}$. It is said to be strictly monotone decreasing if $x_n >
x_{n+1}$ for all $n \in \mathbb{N}$.
\end{defn}

\begin{defn}\label{s1d7}
A sequence $\{x_n\}$ is said to be a Cauchy sequence if for a given $\epsilon
> 0$, we can find $N \in \mathbb{N}$ such that for all $m, n > N$, $|x_m - x_n| <
\epsilon$.
\end{defn}

\begin{defn}\label{s1d8}
If $\{n_1, n_2, \ldots\}$ is a subset of natural numbers such that $n_1 < n_2
< \ldots$ then $\{x_{n_1}, x_{n_2}, \ldots\}$ is said to be a subsequence of 
$\{x_n\}$.
\end{defn}

\begin{defn}\label{s1d9}
If $\{x_n\}$ is a sequence of reals then
\[
\limsup x_n = \inf_{n \ge 1}(\sup_{m \ge n} x_m)
\]
and
\[
\liminf x_n = \sup_{n \ge 1}(\inf_{m \ge n} x_m)
\]
\end{defn}

\begin{prop}\label{s1p1}
If $\{x_n\}$ is convergent then it has a unique limit.
\end{prop}
\begin{proof}
Let if possible, $\lim x_n = L$ and $\lim x_n = M$. Then for a given $\epsilon
> 0$, we can find, $N_1, N_2 \in \mathbb{N}$ such that for all $n > N_1$, 
$|x_n - L| < \epsilon/2$ and for all $n > N_2$, $|x_n - M| < \epsilon/2$. Then, 
if $N = \max(N_1, N_2)$, $|L - M| = |L - x_n + x_n - M| \le |L - x_n| + 
|M - x_n| < \epsilon$.
\end{proof}

\begin{thm}[Monotone Convergence Theorem]\label{s1t1}
If $\{x_n\}$ is a monotone increasing (decreasing) sequence that is bounded
above (below) then $\{x_n\}$ is convergent and its limit is its supremum 
(infimum).
\end{thm}
\begin{proof}
If $\{x_n\}$ is a monotone increasing functions that is bounded above then the
set of the sequence values is a set of real numbers that is bounded above. By 
the completeness property of real numbers, it has a supremum $M \in \mathbb{R}$.
Fix an $\epsilon > 0$. Then there is at least one $x_m$ such that $M - \epsilon
< x_m$. For if there were none such then $M$ would not be the supremum. Since
our sequence is monotone increasing, this means that for all $n > m$, $|M - x_n|
< \epsilon$ so that $M = \lim x_n$.

The proof for monotone decreasing sequence is similar.
\end{proof}

\begin{thm}[Bolzano-Weierstrass Theorem]\label{s1t2}
Every bounded sequence has a convergent subsequence.
\end{thm}
\begin{proof}
Let $\{x_n\}$ be such that $m \le x_n \le M$ for all $n \in \mathbb{N}$. An 
index $n$ of the sequence is called a peak if $x_n > x_m$ for all $m > n$. 
Likewise, an index $n$ of the sequence is called a trough if $x_n \le x_m$
for all $m > n$. Build a subsequence $\{x_{n_k}\}$ inductively as follows. Let 
$n_1 = 1$ and having chosen $n_k$, let $n_{k+1}$ be the peak among indices 
strictly beyond $n_k$. If we can always get a peak at all stages then 
$\{x_{n_k}\}$ is a monotone decreasing sequence that is bounded from below. 
Therefore, by monotone convergence theorem, it converges to its infimum. If, 
after $x_{n_k}$ we cannot find a peak then among the remaining terms we can find
a subsequence of troughs in the same way. It is a monotone increasing sequence
that is bounded from above and hence convergent to its supremum.
\end{proof}

\begin{prop}\label{s1p2}
Every convergent sequence is a Cauchy sequence.
\end{prop}
\begin{proof}
Let $\{x_n\}$ be such that $\lim x_n = L$. Then for a fixed $\epsilon > 0$, we
can find $N \in \mathbb{N}$ such that for all $n > N$, $|x_n - L| < \epsilon/2$.
Therefore, if $m, n > N$, $|x_m - x_n| = |x_m - L + L - x_n| \le |x_m - L| + 
|x_n - L| < \epsilon$, making $\{x_n\}$ a Cauchy sequence.
\end{proof}

\begin{prop}\label{s1p3}
Every Cauchy sequence is bounded.
\end{prop}
\begin{proof}
Let $\{x_n\}$ be a Cauchy sequence. Fix an $\epsilon > 0$. There exists $N \in
\mathbb{N}$ such that for all $m, n \ge N$, $|x_m - x_n| < \epsilon$. In 
particular, for all $m > N$, $|x_m - x_N| < \epsilon$ so that for the tail of
the sequence beyond $x_N$, $x_N - \epsilon$ is the lower bound and $x_N +
\epsilon$ is the upper bound. For the finite set $\{x_1, \ldots, x_{N-1}\}$, let
$m$ and $M$ be the minimum and the maximum. Then $\min(m, x_N - \epsilon)$ is
a lower bound of the entire sequence and $\max(M, x_N + \epsilon)$ is its upper
bound, making $\{x_n\}$ a bounded sequence.
\end{proof}

\begin{prop}\label{s1p4}
Every Cauchy sequence is convergent.
\end{prop}
\begin{proof}
By proposition \ref{s1p3} a Cauchy sequence is bounded. Using Bolzano-Weierstrass
theorem, \ref{s1t2}, it has a convergent subsequence $\{x_{n_k}\}$. Let $L$ be
the limit of this subsequence. For a given $\epsilon > 0$, we can find $N_1 \in
\mathbb{N}$ such that for all $n_k > N_1$, $|x_{n_k} - L| < \epsilon/2$. For the
same $\epsilon$, since $\{x_n\}$ is a Cauchy sequence, there exists $N_2 \in
\mathbb{N}$ such that for all $m, n > N_2$, $|x_m - x_n| < \epsilon/2$. Let
$N = \max(N_1, N_2)$. Then for $m > N$ and $k$ such that $n_k > N$, $|x_m - L| 
= |x_m - x_{n_k} + x_{n_k} - L| \le |x_m - x_{n_k}| + |x_{n_k} - L| < \epsilon$.
In the last inequality $|x_m - x_{n_k}| < \epsilon/2$ by the Cauchy property
and $|x_{n_k} - L| < \epsilon/2$ because $L$ is the limit of the convergent
subsequence.
\end{proof}
\end{document}
